\begin{table}[H]
\centering
\caption{\textbf{Planning, TOPAS transport, and repeat-run uncertainty settings used in the final risk-analysis branch.}}
\label{tab:planning_transport_parameters}
\small
\begin{tabular}{p{5.2cm}p{7.5cm}}
\toprule
Parameter & Setting \\
\midrule
Peripheral lattice target & $\mathrm{CTV}_{\mathrm{boost}} = \mathrm{GTV} + 5$~mm \\
Vertex prescription concept & 15~Gy per vertex \\
Peripheral prescription concept & 3.0--3.5~Gy to $\mathrm{CTV}_{\mathrm{boost}}$ \\
Typical vertex diameter & 0.8--1.2~cm; 1.0~cm default \\
Candidate support radius & 6~mm in fixed safe-core library \\
GTV contraction for safe-core sampling & 5~mm \\
Minimum OAR clearance for candidate selection & 15~mm \\
Candidate sampling step & 6~mm \\
Minimum inter-vertex spacing & 18~mm \\
Layer spacing where layered patterns used & 30~mm \\
Site-specific cohort size & 10 synthetic H\&N cases \\
Estimated cohort GTV range & 66--326~cc \\
Accepted vertex counts across cohort & 2--5 per case after anatomy-aware pruning \\
TOPAS role & Physical transport engine; not a clinical inverse planner \\
Physical delivery model & Two-component photon lattice surrogate: broad base + focused vertex component \\
Dose combination & Linear weighted sum of base and vertex dose grids \\
Physical calibration targets & PTV $D_{95} \approx 3.5$~Gy and peak mean dose $\approx 15$~Gy \\
Default base-component histories & $1\times10^{6}$ \\
Default vertex-component histories & $2\times10^{6}$ \\
Dose post-processing & Body-aware Gaussian denoising before component-weight calibration \\
Repeat-run uncertainty subset & Cases 02--07 \\
Repeat-run seeds & 11, 22, 33 \\
Repeat-run history scales & 0.5, 1.0, 2.0 \\
\bottomrule
\end{tabular}
\end{table}
